\documentclass[12pt]{article}

\usepackage[utf8]{inputenc}
\usepackage[T1]{fontenc}
\usepackage[polish]{babel}


\title{Dokumentacja funkcjonalna}
\author{Mateusz Kamieniecki \\ Jan Rękas}
\date{\today}


\begin{document}

\maketitle

\section{Cel projektu}
Celem projektu jest stworzenie aplikacji w języku C, umożliwiającej wyznaczanie
współrzędnych wierzchołków grafów planarnych w sposób czytelny i intuicyjny dla
użytkownika. Program pozwala na eksperymentowanie z dwoma różnymi algorytmami
rozmieszczania wierzchołków oraz generowanie wyników w formacie binarnym,
przeznaczonym do dalszego przetwarzania, lub w formacie tekstowym, przyjaznym
dla użytkownika.

\section{Opis funkcjonalności}

Program udostępnia następujące funkcje:

\begin{itemize}
	\item wczytywanie grafu z pliku tekstowego,
	\item wybór algorytmu wyznaczania współrzędnych wierzchołków,
	\item obliczanie współrzędnych wierzchołków zgodnie z wybranym algorytmem,
	\item zapis wyników w formacie binarnym lub tekstowym,
	\item wyświetlanie komunikatów o błędach w przypadku nieprawidłowego formatu danych wejściowych,
\end{itemize}




\end{document}

